\scanpage{200}
\mspage{116}
% archive sheet number: 199
% element: text | top, line 1 | ici un énoncé rectifié
\unsure{ici} un énoncé rectifié.

% element: text | upper third, lines 2-9 | Proposition 4.5.5
\uline{Proposition} \unsure{4.5.5|4.5.6} \textcircled{$1^\circ$} Dans les hypothèses générales de 4.5.2, \deleted{\unsure{des card}} $C,\pi_0,\pi_1,\pi$, les conditions $\unsure{a'|a''},b,c'$ de \deleted{loc. cit.} cette proposition sont équivalentes \inserted{[moyennant $\pi\geq\pi_1$]}, et impliquent la condition ($F$ est $\pi$-accessible). Quand $\pi^{\pi_0}=\pi$ (notamment quand $\pi=2^{\unsure{c|\sigma}}$, avec $c\geq\pi_0$), alors toutes les conditions \unsure{a) à i')} de 4.5.2 sont équivalentes.

% slanted
% element: margin note | diagonal left margin beside heading and first paragraph | NB Pour cette équivalence
\margin{\uline{NB} Pour cette équivalence, l'hyp que $C$ soit stable par $\varinjlim$ de card. $\leq\pi_0$ est inutile, et aussi \deleted{\unsure{(quand} $\pi^{\pi_0}=\pi$)} \deleted{\unsure{d'après le lemme}} \unsure{l'implication} \unsure{c)}, et l'équival. si $\pi=\pi^{\pi_0}\geq\pi_1$, \unsure{cf.} \unsure{4.7.1} (p.\unsure{161|151}). \unsure{(a'|c')} \unsure{et} \unsure{c)}\unsure{?}}

% element: text | middle, lines 9-12 | Plus généralement et plus précisément
\textcircled{$2^\circ$} Plus généralement et plus précisément, \deleted{\unsure{avec}} pour $\pi\geq\pi_0$ donné, \deleted{\unsure{les relations}} considérons les conditions sur $\pi$
\begin{enumerate}[label=\arabic*)]
% element: display | middle, item 1 | pi power pi0 equals pi
\item $\pi^{\pi_0}=\pi$
% element: text | middle, item 2 | Filt pi M
\item $\mathrm{Filt}^{\pi}(M)$ (défini dans 4.5.3) est stable par $\varinjlim$ de cardinal $\leq\pi$
% element: text | lower middle, item 3 | for all F in M
\item $\forall F$ dans $M$, on a $F\simeq\varinjlim_I X_i$, avec $I$ ens. ordonné grand devant $\pi$, les $X_i$ dans $\mathrm{Filt}^{\pi}(M)$. De plus, si $\pi'>\pi$ et $F\in\mathrm{Filt}^{\pi'}(M)$, on peut prendre $I$ de cardinal $\leq\unsure{\pi^{\pi^*}|\pi^{\pi'}|\pi'{}^{\pi}|\pi'{}^{\pi_0}}$.
% element: text | lower third, item 4 | Comme 3
\item Comme 3), sans la majoration de cardinaux, et \deleted{\unsure{avec}} en supposant seulement $I$ catégorie \unsure{grd|grde|grand} devant $\pi$ (pas néc. cat. ordonnée)
% element: text | foot, item 5 | La condition c
\item La condition c) de 4.5.2 implique les autres conditions (elles sont équivalentes, par $1^\circ$), en d'autres termes
% element: display | foot | M pi subset Filt pi
\[
M_\pi\subset\mathrm{Filt}^{\pi}\qquad(\text{donc }M_\pi=\mathrm{Filt}^{\pi}\text{ puisque}
\]
\unsure{on sait} que $\mathrm{Filt}^{\pi}\subset M_\pi$)
\end{enumerate}
